\section{CZARNE WROTA}

Ogólna zasada:

Odnalezienie i zniszczenie Czarnych Wrót jest celem, który oddział musi wypełnić aby wygrać.
Czarne Wrota strzeżone są przez Demony oraz przez Bezimiennego.
Czarne Wrota posiadają w czasie walki własną fazę, podczas której atakują czarem ,,wybuch''.

Sposób postępowania:

Po odnalezieniu Czarnych Wrót należy rzucić 2K3+1, aby ustalić liczbę strzegących je Demonów, a następnie ustawić potwory w szyku, przy czym Demony zajmują pierwsze szeregi, Bezimienny ustawiany jest w ostatnim szeregu Demonów lub za nimi, zaś Czarne Wrota zupełnie z tyłu (w trzecim lub nawet czwartym szeregu).

Podczas swojej fazy Czarne Wrota atakują czarem ,,wybuch'' bez względu na szereg, w którym się znajdują, rzucając go trzykrotnie w ciągu każdej fazy.
Czarne Wrota nie ponoszą kosztu rzucania czaru, przy czym wszystkie trzy czary, rzucone w ciągu jednej fazy mogą być skierowane na tę samą postać lub podzielone między kilka postaci.

Czarne Wrota mają zdolność rzucania czarów dopóty, dopóki żyje Bezimienny.
Po jego śmierci Czarne Wrota tracą swą moc, a wszystkie Demony znikają.
Po zabiciu Bezimiennego i Demonów może zostać podjęta próba zniszczenia Czarnych Wrót.
Musi być ona podjęta przez postać, posiadającą odpowiednie zdolności.
Próba polega na rzuceniu 1K6 i porównaniu wyrzuconej liczby oczek z wartością zdolności ,,Czarne Wrota'' zaangażowanej w nią postaci.
Jeżeli liczba oczek jest mniejsza od zdolności lub im równa próba została zakończona sukcesem – Czarne Wrota rozpadły się w pył.
Niepowodzenie próby oznacza odniesienie przez przeprowadzającą ją postać 1 rany.
Próby mogą być powtarzane aż do
a. zakończenia ich sukcesem lub
b. śmierci wszystkich postaci, zdolnych do ich podjęcia lub
c. rezygnacji i wycofania się oddziału.

W przypadku c. oddział w czasie następnych przygód znajdzie Czarne Wrota w tym samym miejscu, jednak znów strzeżone przez Demony i Bezimiennego.

Z Bezimiennym i Demonami można tylko walczyć – nie mają tu zastosowania przepisy, dotyczące negocjacji i wykupu.
Bezimienny jest zawsze tylko jeden, liczba Demonów jest ustalana wyłącznie rzutem kostki – nie wpływa na nią poziom, na którym Czarne Wrota się znajdują, jakkolwiek zarówno Bezimienny, jak i Demony podlegają pozostałym modyfikacjom, wynikającym z tabeli 6.7 ,,Poziom labiryntu''.
Zdolności ,,Czarne Wrota'' nie mogą przekroczyć 5 punktów.
